如果你真的碰巧忘记备份，而系统又出了难以恢复的问题，或者定位、修复所花的时间已经不如重装一遍，那就可能需要考虑重新安装了。当然，希望各位不要经常用到这里的方法，平时大改环境前养成备份的习惯，会省下很多事。环境变量配错通常还能修回来，不用一报“找不到命令”就准备重装；我们先根据症状排查，必要时用 Live U 盘确认数据还能不能读取，再决定下一步。需要回顾的话，安装方法见第\ref{sec:ubuntu-partition}节，快照和个人文件备份见第\ref{sec:system-backup}节。

\begin{note}
\textbf{格式化会清空所选分区。}先把源码、数据、密钥及所需配置等重要数据复制到另一块介质，并实际读取验证。记录磁盘型号、根分区、已有 ESP 和独立 /home 的身份。若 /home 与根目录在同一分区，格式化根分区也会删除个人文件；即使 /home 独立，也必须在安装器中明确取消它的格式化。保留 ESP 与 Windows 分区，不要直接删除整块盘。
\end{note}

\Needspace{8\baselineskip}
\begin{example}[重装前的核对流程]
重装过程其实还是原先那条安装路线，只是这一次，磁盘里已经有我们的文件了，所以要更仔细地分清哪些保留、哪些替换。

首先进入 Live 试用环境，运行 \texttt{lsblk -o NAME,SIZE,FSTYPE,MOUNTPOINTS,MODEL}，把输出和第一次安装时的记录对应起来。然后打开备份介质，实际检查一个源码文件和一个数据文件，确认备份不是只留下了几个文件名。

接下来按照原先的安装方法，一直运行到“其它选项”这一步。只对计划替换的 Ubuntu 根分区设置格式化；如果 /home 是独立分区，就重新指定它的挂载点，\textbf{不勾选格式化}，已有 ESP 也复用而不格式化。最后再读一遍写盘摘要，发现分区身份或备份还没弄清楚，就先退出，不要抱着“应该是这个”的心态点下去。

安装完成后，分别启动 Ubuntu 和 Windows，核对挂载点，再恢复自己的文件，重新构建一个工程。能进入桌面只是第一步，文件还在、工程也能继续用，才算把工作环境接回来了。

如果只是学习这里的方法，不必为了练习把主力机重装一遍。你可以在虚拟机或空白练习盘上完成一次重装并记录结果。
\end{example}
\begin{exercise}\label{ex:v1-reinstall-plan}
假设根分区损坏、独立 /home 仍可读取，写一份重装清单：备份验证、分区身份、格式化范围、ESP 处理和安装后的三个核验动作。再说明 /home 没有独立分区时要怎样调整方案。
\end{exercise}
